\chapter{1956年全国卷}

\begin{problem}
\begin{enumerate}
\item 利用对数性质，计算 \(\lg^{2}5+\lg2\cdot\lg50\). （注）\(\lg\) 是以 \(10\) 为底的对数 \(\log_{10}\) 的记号.

\item 设 \(m\) 是实数，求证方程 \(2x^{2}-\left(4m-1\right)x-m^{2}-m=0\) 的两个根必定都是实数.

\item 设 \(M\) 是 \(\triangle ABC\) 的边 \(AC\) 的中点，过 \(M\) 作直线交 \(AB\) 直线于 \(E\)，过 \(B\) 作直线平行于 \(ME\) 交 \(AC\) 直线于 \(F\). 求证 \(\triangle AEF\) 的面积等于 \(\triangle ABC\) 的面积的一半.

\item 一个三角形三边的长分别是 \(3\text{ 尺}\)，\(4\text{ 尺}\) 及 \(\sqrt{37}\text{ 尺}\). 求这个三角形的最大角的度数.

\item 设 \(\tan\alpha\) 与 \(\tan\beta\) 是方程 \(x^{2}+6x+7=0\) 的两个根，求证 \(\sin\left(\alpha+\beta\right)=\cos\left(\alpha+\beta\right)\).
\end{enumerate}
\end{problem}

\begin{solution}
\begin{enumerate}
\item 由 \(\lg\left(2\cdot5\right)=\lg10=1\)，得 \(\lg2+\lg5=1\)，所以 \(\lg2=1-\lg5\). 又因为 \(50=10\cdot5\)，所以 \(\lg50=\lg10+\lg5=1+\lg5\). 代入原式，得
\[
\lg^{2}5+\lg2\cdot\lg50=\lg^{2}5+\left(1-\lg5\right)\left(1+\lg5\right)=1
\]
因此原式的值为 \(1\).

\item 方程的二次项系数为 \(2\)，不为零. 其判别式为
\[
\Delta=\left(4m-1\right)^{2}-4\cdot2\cdot\left[-\left(m^{2}+m\right)\right]=\left(4m-1\right)^{2}+8\left(m^{2}+m\right)=24m^{2}+1
\]
因为 \(m\) 是实数，所以 \(m^{2}\geq0\)，从而 \(\Delta=24m^{2}+1>0\). 因此该方程有两个互不相等的实数根，结论得证.

\item 取直角坐标系，并以 \(AC/2\) 为长度单位，使 \(A=(0,0)\)，\(C=(2,0)\)，于是 \(M=(1,0)\). 设 \(B=(u,v)\)，其中 \(v\ne0\). 因为 \(E\) 在直线 \(AB\) 上，可设 \(E=(tu,tv)\). 这里 \(t\ne0\)，否则 \(E=A\)，直线 \(ME\) 与 \(AC\) 重合，过 \(B\) 作的平行线不与 \(AC\) 相交，构造不成立.

由 \(M\)，\(E\) 两点坐标，得 \(\overrightarrow{ME}=(tu-1,tv)\). 直线 \(BF\) 与 \(ME\) 平行，所以直线 \(BF\) 上的点可表示为
\[
B+\lambda\overrightarrow{ME}=\left(u+\lambda\left(tu-1\right),v+\lambda tv\right)
\]
由于 \(F\) 在直线 \(AC\) 上，其纵坐标为零，故 \(v+\lambda tv=0\). 因为 \(v\ne0\) 且 \(t\ne0\)，所以 \(\lambda=-t^{-1}\)，从而
\[
F=\left(u-\frac{tu-1}{t},0\right)=\left(t^{-1},0\right)
\]
于是
\[
S_{\triangle AEF}=\frac{1}{2}\left|\det\begin{pmatrix}tu&tv\\ t^{-1}&0\end{pmatrix}\right|=\frac{|v|}{2}
\]
另一方面，\(C=(2,0)\)，所以
\[
S_{\triangle ABC}=\frac{1}{2}\left|\det\begin{pmatrix}u&v\\ 2&0\end{pmatrix}\right|=|v|
\]
因此 \(S_{\triangle AEF}=\frac{1}{2}S_{\triangle ABC}\)，结论得证.

\item 因为 \(\sqrt{37}>4>3\)，所以最大角 \(\theta\) 的对边为 \(\sqrt{37}\). 由余弦定理，
\[
37=3^{2}+4^{2}-2\cdot3\cdot4\cos\theta
\]
即 \(\cos\theta=-\frac{1}{2}\). 三角形内角满足 \(0^{\circ}<\theta<180^{\circ}\)，故 \(\theta=120^{\circ}\).

\item 由韦达定理，\(\tan\alpha+\tan\beta=-6\)，\(\tan\alpha\cdot\tan\beta=7\). 题设说明 \(\tan\alpha\)、\(\tan\beta\) 有定义，且 \(1-\tan\alpha\tan\beta=1-7=-6\ne0\)，所以正切和角公式适用. 因而
\[
\tan\left(\alpha+\beta\right)=\frac{\tan\alpha+\tan\beta}{1-\tan\alpha\tan\beta}=\frac{-6}{-6}=1
\]
又因 \(\cos\left(\alpha+\beta\right)\ne0\)，故
\[
\frac{\sin\left(\alpha+\beta\right)}{\cos\left(\alpha+\beta\right)}=1
\]
从而 \(\sin\left(\alpha+\beta\right)=\cos\left(\alpha+\beta\right)\)，结论得证.
\end{enumerate}
\end{solution}

\begin{problem}
解联立方程 \(\begin{cases}
7\sqrt{x+y}-x-y=12,\\
x^{2}+y^{2}=136.
\end{cases}\)
\end{problem}

\begin{solution}
由根式有意义知 \(x+y\geq0\). 令 \(u=\sqrt{x+y}\)，则 \(u\geq0\)，且第一个方程化为 \(7u-u^{2}=12\)，即
\(u^{2}-7u+12=0\)，\(\left(u-3\right)\left(u-4\right)=0\)
所以 \(u=3\) 或 \(u=4\)，从而 \(x+y=9\) 或 \(x+y=16\).

又因为
\[
2\left(x^{2}+y^{2}\right)-\left(x+y\right)^{2}=x^{2}-2xy+y^{2}=\left(x-y\right)^{2}
\]
所以可在两种情形下分别求出 \(x-y\)，再利用 \(x=\frac{(x+y)+(x-y)}{2}\)，\(y=\frac{(x+y)-(x-y)}{2}\) 求出 \(x\)，\(y\).

当 \(x+y=9\) 时，
\[
\left(x-y\right)^{2}=2\cdot136-9^{2}=191
\]
所以 \(x-y=\pm\sqrt{191}\). 当 \(x-y=\sqrt{191}\) 时，
\(x=\frac{9+\sqrt{191}}{2}\)，\(y=\frac{9-\sqrt{191}}{2}\)
当 \(x-y=-\sqrt{191}\) 时，
\(x=\frac{9-\sqrt{191}}{2}\)，\(y=\frac{9+\sqrt{191}}{2}\)

当 \(x+y=16\) 时，
\[
\left(x-y\right)^{2}=2\cdot136-16^{2}=16
\]
所以 \(x-y=\pm4\). 当 \(x-y=4\) 时，\(x=10\)，\(y=6\)；当 \(x-y=-4\) 时，\(x=6\)，\(y=10\).

这四组数都满足 \(x+y\geq0\)，并且由构造满足 \(x+y=9\) 或 \(x+y=16\) 以及 \(x^{2}+y^{2}=136\). 代回第一个方程，分别得到 \(7\cdot3-9=12\) 或 \(7\cdot4-16=12\). 因此原联立方程的解为
\[
\left\{
\left(\frac{9+\sqrt{191}}{2},\frac{9-\sqrt{191}}{2}\right),
\left(\frac{9-\sqrt{191}}{2},\frac{9+\sqrt{191}}{2}\right),
(10,6),(6,10)
\right\}
\]
\end{solution}

\begin{problem}
设 \(P\) 是等边三角形 \(ABC\) 外接圆弧 \(\myarc{BC}\)（不含 \(A\)）上的一点，求证 \(PA^{2}=AB^{2}+PB\cdot PC\).
\end{problem}

\begin{solution}
设直线 \(AP\) 与 \(BC\) 相交于 \(M\). 因为 \(P\) 在弧 \(BC\)（不含 \(A\)）上，\(M\) 在线段 \(BC\) 上，且 \(AM+MP=AP\).

\begingroup
\FigureLayoutDeclare{img/1956/national/q3_solution_fig1.png}{5.8cm}{0bp 0bp 0bp 0bp}
\solutionfigure{\bitmapfigure[max width=4.4cm]{img/1956/national/q3_solution_fig1.png}}
\endgroup

由于 \(A\)，\(B\)，\(P\)，\(C\) 四点共圆，\(\angle APB=\angle ACB=\angle ABM\)，又因 \(A\)，\(P\)，\(M\) 三点共线，\(\angle PAB=\angle MAB\). 所以 \(\triangle APB\sim\triangle AMB\)，从而
\[
AP\cdot AM=AB^{2}
\]

等边三角形的内角均为 \(60^{\circ}\). 由圆周角定理，\(\angle APC=\angle BPA=60^{\circ}\). 因为 \(P\)，\(M\)，\(A\) 三点共线，所以 \(\angle APC=\angle BPM\)；又因为 \(A\)，\(B\)，\(P\)，\(C\) 四点共圆，\(\angle PAC=\angle PBC\)，而 \(B\)，\(M\)，\(C\) 三点共线，故 \(\angle PAC=\angle PBM\). 因此 \(\triangle APC\sim\triangle BPM\)，从而
\[
AP\cdot MP=PB\cdot PC
\]

将两式相加，得 \(AP\cdot AM+AP\cdot MP=AB^{2}+PB\cdot PC\). 由 \(AM+MP=AP\)，左边等于 \(AP^{2}\)，所以 \(PA^{2}=AB^{2}+PB\cdot PC\)，结论得证.
\end{solution}

\begin{problem}
有一四棱柱体，底面 \(ABCD\) 为菱形，\(\angle A'AB=\angle A'AD\)（如图），求证平面 \(A'ACC'\) 垂直于底面 \(ABCD\).

\bitmapfigure[max width=4.8cm]{img/1956/national/q4_fig1.png}
\end{problem}

\begin{solution}
设 \(AC\) 与 \(BD\) 交于 \(O\). 因为 \(ABCD\) 是菱形，所以两条对角线互相垂直且互相平分，故 \(AO\perp BD\)，且 \(O\) 是 \(BD\) 的中点.

\begingroup
\FigureLayoutDeclare{img/1956/national/q4_solution_fig1.png}{5.8cm}{0bp 0bp 0bp 0bp}
\solutionfigure{\bitmapfigure[max width=4.4cm]{img/1956/national/q4_solution_fig1.png}}
\endgroup

因为 \(AB=AD\)，\(A'A\) 是公共边，且 \(\angle A'AB=\angle A'AD\)，所以由边角边可知 \(\triangle A'AB\cong\triangle A'AD\)，从而 \(A'B=A'D\). 因此 \(\triangle A'BD\) 是等腰三角形. 又因为 \(O\) 是 \(BD\) 的中点，所以等腰三角形的中线 \(A'O\) 同时是高，即 \(A'O\perp BD\).

直线 \(AO\) 与 \(A'O\) 都在平面 \(A'ACC'\) 内，且相交于 \(O\)，并且都垂直于直线 \(BD\)，所以由直线与平面垂直的判定定理，\(BD\) 垂直于平面 \(A'ACC'\). 由于 \(BD\) 在底面 \(ABCD\) 内，且底面 \(ABCD\) 含有直线 \(BD\)，由平面与平面垂直的判定定理，平面 \(A'ACC'\) 垂直于底面 \(ABCD\)，结论得证.
\end{solution}

\begin{problem}
若三角形的三个角成等差级数，则其中一定有一个角是 \(60^{\circ}\)；若这样的三角形的三边又成等比级数，则三个角都是 \(60^{\circ}\)，试证明之.
\end{problem}

\begin{solution}
将三角形的三个角按从大到小记为 \(A\)，\(B\)，\(C\). 因为它们成等差数列，所以 \(A+C=2B\). 又因为三角形内角和为 \(A+B+C=180^{\circ}\)，两式相减得 \(3B=180^{\circ}\)，所以 \(B=60^{\circ}\). 这证明了三个角中一定有一个是 \(60^{\circ}\).

再按从小到大排列三边，设它们为 \(a\)，\(aq\)，\(aq^{2}\)，其中 \(a>0\)，\(q\geq1\). 三角形中边长大小规律与其对角大小规律相同，因此中间边 \(aq\) 所对的角正是中间角 \(B=60^{\circ}\). 对该角应用余弦定理，得
\[
(aq)^{2}=a^{2}+(aq^{2})^{2}-2a\cdot aq^{2}\cos60^{\circ}=a^{2}+a^{2}q^{4}-a^{2}q^{2}
\]
由于 \(a>0\)，两边除以 \(a^{2}\)，得 \(q^{2}=1+q^{4}-q^{2}\)，即
\(q^{4}-2q^{2}+1=0\)，\(\left(q^{2}-1\right)^{2}=0\)
所以 \(q^{2}=1\). 结合 \(q\geq1\)，得 \(q=1\)，三边相等. 因此三角形为等边三角形，三个角都是 \(60^{\circ}\)，结论得证.
\end{solution}
